QtRocket is an open-source model-rocket simulator: a C++23 simulation core shared by a
Qt6 Widgets GUI, a Qt-free headless REPL, and a standalone OpenGL design viewer. This
whitepaper is a whole-system architecture review of the working tree at commit
\code{254cd41} --- what the program can do today, how its layers fit together, and where
it stands relative to OpenRocket, used throughout as a comparison baseline for an
independent project with similar goals, not as a specification to clone.

What exists is a genuinely finished, well-guarded 3-DOF point-mass simulator: a
composite part tree whose placement is stored as physical intent and resolved by a
single geometry resolver shared with the visualizer, time-aware composite mass/CG/inertia
behind a mass-delta cache, motor ingest from RockSim files, RASP files, and a live
thrustcurve.org client, a validated seven-layer US Standard Atmosphere 1976, pluggable
gravity models, fixed-step RK4 and adaptive RKF45 integration, typed termination
diagnostics around a hang-proof run loop, XML design persistence, and a 33-command
scriptable REPL that exercises all of it. Around that core sits the review's central
finding: a ring of machinery that is built, unit-tested, and consumed by nothing. The
Barrowman center-of-pressure composition, the per-step inertia tensor (recorded every step, integrated by nothing), the wind model,
the quaternion state slots, and the torque interface all exist and all have zero
production consumers. Flight physics therefore remains thrust-along-world-vertical plus
gravity plus a single constant-$C_d$ drag term: no angle of attack, no restoring moment,
no weathercocking, no recovery deployment --- every flight ends ballistically. The GUI
lags furthest behind the core: rocket design is possible only from the CLI, with the
design-editor widgets empty shells at the reviewed commit and in-flight work on a
separate branch.

The review was produced by eight independent subsystem deep-reads, each adversarially
verified citation-by-citation against the source, plus live execution of the CLI binary;
design documents under \srcf{docs/} were treated as claims to test, not facts to
inherit. The result is a feature matrix against OpenRocket, a ranked inventory of
functional deficiencies --- flight-dynamics fidelity first, recovery second, component
catalog and design UI close behind --- and a suggested sequencing that consumes the
already-built seams before building new ones.

\vspace{0.5em}
\begin{designrule}[The review in one sentence]
QtRocket's architecture is ahead of its features: the discipline --- fail-closed
diagnostics, strict warnings, five-platform CI, verified seams --- is that of a much
larger simulator, and the shortest path to realistic flights is to \emph{consume} the
dormant machinery (Barrowman aerodynamics, torques, wind, recovery events), not to
build new subsystems.
\end{designrule}
